% Highlights (max 5 bullets, <=85 caracteres cada uno incluyendo
% espacios — norma de Elsevier). Se suben como fichero aparte
% ("Highlights") en el Editorial Manager.
%
% Revision 2026-09: el cuarto punto decia que las reglas muestreadas
% "igualan" al genetico publicado con mejor-de-1024. El articulo mide
% 11.3 contra 9.8, mejor en 3 de 12, asi que "igualan" no se sostiene y
% es justo la implicacion de eficiencia que la revision r2 senala. El
% quinto decia que los schedules son mas robustos "que las baselines",
% sin decir que eso exige el objetivo robusto. Los dos reescritos.
\documentclass{article}
\usepackage[T1]{fontenc}
\usepackage{lmodern}
\begin{document}
\section*{Highlights}
\begin{itemize}
\item First GP hyper-heuristic for the job shop scheduling problem with
interval durations
\item Evolved rules beat every constructive baseline on all 70 interval
Taillard instances
\item Rules evolved on four $20\times15$ instances generalize zero-shot
across sizes and families
\item Interval widths pay under a robustness objective, not under
expected makespan
\item A genetic algorithm needs 76,000 schedules to match one pass of the
evolved rule
\end{itemize}
\end{document}
